\appendix
\clearpage
\section{Derivations}
\label{app:math}

Everything used in the body is derived here, and every identity carries a tag
(\textsf{A.1}, \textsf{B.5}, \dots) that is also the name of a numerical check
in \texttt{verify\_derivations.py}. A derivation that is only written down can
be wrong in a way no reader notices; all $47$ of them are therefore evaluated
as well, and Table~\ref{tab:derivs} lists the residual of each.

\subsection{Conventions and the two-site algebra}

\paragraph{\textsf{A.1} Pauli versus spin.}
With $\bm\sigma = 2\bm S$,
\begin{equation}
  \ssig_i\cdot\ssig_j = 4\,\bm S_i\cdot\bm S_j .
  \tag{A.1}
\end{equation}
Every textbook result quoted below is stated for $\bm S\cdot\bm S$ and carries
this factor of $4$ when translated to the Hamiltonians as written in the
assignment. It is the single easiest way to be wrong by $4$ in this problem.

\paragraph{\textsf{A.2}--\textsf{A.3} Spectrum of the exchange operator.}
Let $P_{ij}$ be the operator that swaps sites $i$ and $j$. A one-line
computation on the four basis states gives
\begin{equation}
  P_{ij} = \tfrac12\big(\mathbb{1} + \ssig_i\cdot\ssig_j\big)
  \quad\Longleftrightarrow\quad
  \ssig_i\cdot\ssig_j = 2P_{ij} - \mathbb{1}.
  \tag{A.3}
\end{equation}
Since $P^2=\mathbb{1}$, $P$ has eigenvalues $\pm1$: $+1$ on the symmetric
(triplet) subspace and $-1$ on the antisymmetric (singlet) one. Hence
\begin{equation}
  \operatorname{spec}\big(\ssig\cdot\ssig\big) = \{\,+1 \text{ (triplet, 3-fold)},\;
  -3 \text{ (singlet)}\,\}.
  \tag{A.2}
\end{equation}
This single fact drives everything in Model 1: the gate, its cost, its
symmetry, and its differentiation rule.

\paragraph{\textsf{A.6} The gate in closed form.}
On $\{\ket{00},\ket{11}\}$, $\ssig\cdot\ssig$ acts as the identity; on
$\{\ket{01},\ket{10}\}$ its matrix is $\begin{psmallmatrix}-1&2\\2&-1\end{psmallmatrix}
= -\mathbb{1} + 2X$. Exponentiating each block,
\begin{equation}
  e^{-i\theta\,\ssig\cdot\ssig} =
  \begin{cases}
    e^{-i\theta} & \text{on } \ket{00},\ket{11},\\[2pt]
    e^{+i\theta}\big(\cos 2\theta\,\mathbb{1} - i\sin 2\theta\,X\big)
      & \text{on } \ket{01},\ket{10}.
  \end{cases}
  \tag{A.6}
\end{equation}
Applying this by slicing a reshaped state vector, rather than contracting a
generic $4\times4$ tensor, is what makes $L=16$ with $P=10$ layers run in
minutes rather than hours.

\paragraph{\textsf{A.4}--\textsf{A.5} Circuit form.}
On a fixed pair, $(XX)(YY) = (XY)\otimes(XY) = (iZ)\otimes(iZ) = -ZZ$ and
$(YY)(XX) = (-iZ)\otimes(-iZ) = -ZZ$ as well, so $XX$ and $YY$ commute; the
other two pairs are identical by symmetry. Commuting exponents factorise, so
\begin{equation}
  e^{-i\theta(XX+YY+ZZ)} = R_{XX}(2\theta)R_{YY}(2\theta)R_{ZZ}(2\theta),
  \qquad R_{PP}(\varphi) \equiv e^{-i\varphi PP/2}.
  \tag{A.5}
\end{equation}

\subsection{Model 1: ansatz, symmetry and gradients}

\paragraph{\textsf{B.0} The Bethe benchmark.}
For $H_S = J\sum \bm S_i\cdot\bm S_{i+1}$ the Bethe ansatz gives
$E/L \to J(\tfrac14 - \ln 2)$. With (A.1),
\begin{equation}
  \frac{E_0}{L} \longrightarrow 4J\left(\tfrac14-\ln 2\right)
  = J(1-4\ln 2) = -1.7725887\,J .
  \tag{B.0}
\end{equation}

\paragraph{No Trotter error.}
For even $L$ with periodic boundaries the bonds split into
$B_{\rm odd}=\{(0,1),(2,3),\dots\}$ and
$B_{\rm even}=\{(1,2),\dots,(L{-}1,0)\}$, each a set of \emph{disjoint} pairs.
Disjoint two-site operators commute, so
$e^{-i\theta H_{\rm odd}} = \prod_{(a,b)\in B_{\rm odd}} e^{-i\theta\,\ssig_a\cdot\ssig_b}$
\emph{exactly}. The circuit implements the ansatz with no splitting error, so
the only error in the tables is variational.

\paragraph{\textsf{B.1}--\textsf{B.3} Why the ansatz cannot leave the physical
sector.}
Each generator $\ssig_a\cdot\ssig_b$ is an $SU(2)$ scalar: writing
$\bm S_{\rm tot} = \sum_i \bm S_i$, one has
$[\bm S_{\rm tot}, \ssig_a\cdot\ssig_b]=0$ because the dot product of two
vector operators is invariant. Hence every gate --- and therefore the whole
circuit --- commutes with $\bm S_{\rm tot}^2$ and $S^z_{\rm tot}$. The
reference state is a product of two-site singlets, for which
$\bm S^2_{\rm tot}\ket{\psi_0}=0$ and $S^z_{\rm tot}\ket{\psi_0}=0$, so
\begin{equation}
  \bm S_{\rm tot}^2\ket{\psi(\bm\theta)} = 0,\qquad
  S^z_{\rm tot}\ket{\psi(\bm\theta)} = 0
  \quad\text{for all }\bm\theta .
  \tag{B.2--B.3}
\end{equation}
The antiferromagnetic ground state of the even-$L$ chain lives in exactly that
sector (Marshall's theorem), so no expressivity is wasted on states that are
orthogonal to the target by symmetry. Numerically, the amplitude of
$\ket{\psi(\bm\theta)}$ outside $S^z=0$ is $0$ to machine precision and
$\|\bm S^2_{\rm tot}\ket{\psi}\| \sim 10^{-15}$, at random $\bm\theta$.

\paragraph{\textsf{B.4} Adjoint gradients.}
Write $\ket{\psi} = U_M \cdots U_1\ket{\psi_0}$ with $U_k = e^{-i\theta_k H_k}$
and $\ket{\psi_k} = U_{k-1}\cdots U_1\ket{\psi_0}$. Since $H_k$ is a sum of
\emph{disjoint} terms, $\partial_{\theta_k}U_k = -iH_kU_k$, and
\begin{align}
  \frac{\partial E}{\partial\theta_k}
  &= 2\,\mathrm{Re}\,\bra{\psi}H\,U_M\cdots(\partial_{\theta_k}U_k)\cdots U_1\ket{\psi_0}
   \nonumber\\
  &= 2\,\mathrm{Re}\,\bra{\lambda_k}(-iH_k)\ket{\psi_{k+1}},
  \qquad
  \ket{\lambda_k} \equiv \Big(\prod_{j>k}U_j\Big)^{\!\dagger} H \ket{\psi}.
  \tag{B.4}
\end{align}
Both $\ket{\psi_{k+1}}$ and $\ket{\lambda_k}$ are obtained by sweeping once
forwards and once backwards, so the whole gradient costs about two forward
passes, independent of the number of parameters. Measured against central
differences: agreement $\le 3\times10^{-9}$, speed-up $\times 8$ at $L=16$.

\paragraph{\textsf{B.5}--\textsf{B.6} Parameter-shift rules, correctly.}
Let $U(\theta)=e^{-i\theta G}$. Expanding in the eigenbasis of $G$, the matrix
elements of $U^\dagger H U$ carry phases $e^{i(\lambda_j-\lambda_k)\theta}$, so
$E(\theta)$ is a trigonometric polynomial whose frequencies are the
\emph{distinct positive gaps} of $\operatorname{spec}(G)$.

For a \emph{single bond}, (A.2) gives gaps $\{0,4\}$: a single frequency
$\Omega=4$, hence $E(\theta)=a+b\cos4\theta+c\sin4\theta$ and the two-term rule
applies in rescaled form,
\begin{equation}
  \frac{dE}{d\theta}
  = \frac{\Omega}{2}\left[E\!\left(\theta+\frac{\pi}{2\Omega}\right)
    - E\!\left(\theta-\frac{\pi}{2\Omega}\right)\right]
  = 2\left[E\!\left(\theta+\tfrac{\pi}{8}\right)-E\!\left(\theta-\tfrac{\pi}{8}\right)\right].
  \tag{B.5}
\end{equation}
The natural guess is wrong in both directions: the \emph{naive} $\pi/2$-shift,
$\tfrac12$-prefactor version fails, while the rescaled one above is exact ---
verified to $3\times10^{-10}$.

For a whole sub-layer, $G=\sum_{r=1}^{R}(\ssig\cdot\ssig)_r$ over $R=L/2$
disjoint bonds. Its eigenvalues are $R-4k$, $k=0,\dots,R$, so the positive gaps
are $4,8,\dots,4R$: $R$ distinct frequencies, and an exact rule needs $2R$
evaluations. Numerically, at $L=8$ ($R=4$) the two-term rule is off by
$2.2$ in absolute terms --- not a small error, a wrong answer --- while the
$2R$-term rule is exact to $1.4\times10^{-9}$. With $2P$ parameters this would
cost $2P\cdot 2R = 2PL$ circuit evaluations per gradient against two passes
for (B.4), which is why the code uses (B.4).

\paragraph{\textsf{B.7} Why the $1/L^2$ extrapolation misses by $3\times10^{-3}$.}
The XXX chain is critical, described by a $c=1$ conformal field theory, for
which $E_0(L)/L = e_\infty - \pi v c/(6L^2) + \dots$. The subleading terms are
\emph{not} a clean series in $1/L^2$: the isotropic point carries a marginally
irrelevant operator that produces corrections of relative order
$1/\ln^3 L$. A two-parameter fit on five sizes therefore cannot land on
$e_\infty$ to better than about $10^{-3}$, and does not. Quoting the deviation
as a failure of the fit rather than as a physical feature would be the mistake.

\subsection{Model 2: Haldane--Shastry}

\paragraph{\textsf{C.1} The geometric identity.}
With $z_j=e^{2\pi i j/L}$,
\begin{align}
  |z_j-z_k|^2
  &= \big|e^{2\pi ij/L}-e^{2\pi ik/L}\big|^2
   = \big|e^{\pi i(j+k)/L}\big|^2\,
     \big|e^{\pi i(j-k)/L}-e^{-\pi i(j-k)/L}\big|^2 \nonumber\\
  &= 1\cdot\big|2i\sin\big(\pi(j-k)/L\big)\big|^2
   = 4\sin^2\!\Big(\frac{\pi(j-k)}{L}\Big).
  \tag{C.1}
\end{align}

\paragraph{\textsf{C.2} The Hamiltonian as written \emph{is} the textbook one.}
Substituting (C.1) and (A.1), the two factors of $4$ cancel:
\begin{equation}
  H_{\rm HS}
  = \frac{\pi^2 J}{L^2}\sum_{j<k}\frac{\ssig_j\cdot\ssig_k}{|z_j-z_k|^2}
  = \frac{\pi^2 J}{L^2}\sum_{j<k}\frac{4\,\bm S_j\cdot\bm S_k}
      {4\sin^2(\pi(j-k)/L)}
  = J\Big(\frac{\pi}{L}\Big)^{2}\sum_{j<k}
      \frac{\bm S_j\cdot\bm S_k}{\sin^2(\pi(j-k)/L)} .
  \tag{C.2}
\end{equation}
The complex-coordinate dressing of the problem statement is therefore
cosmetic --- but only once (C.1) is in hand.

\paragraph{\textsf{C.3} Closed-form energy.}
Given (C.2) the standard Haldane--Shastry result applies verbatim,
\begin{equation}
  E_0 = -\frac{\pi^2 J}{24}\Big(L+\frac{5}{L}\Big),
  \qquad \frac{E_0}{L}\xrightarrow[L\to\infty]{} -\frac{\pi^2}{24}
  = -0.4112335 .
  \tag{C.3}
\end{equation}
Verified against exact diagonalisation at $L=4,6,8,10,12,14$: agreement
$\le 9\times10^{-15}$ at every size. It is used as a benchmark only after that
check, not on authority.

\paragraph{\textsf{C.4}--\textsf{C.6} The exact ground state, and why the
momentum factor matters.}
In the $S^z=0$ sector, with down spins at sites $n_1<\dots<n_M$, $M=L/2$,
\begin{equation}
  \psi_{\rm J}(n_1,\dots,n_M)\;\propto\;
  \prod_{i=1}^{M} z_{n_i}\;\prod_{i<j}\big(z_{n_i}-z_{n_j}\big)^2 .
  \tag{C.4}
\end{equation}
The Jastrow factor $\prod_{i<j}(z_{n_i}-z_{n_j})^2$ is a squared Vandermonde
determinant --- the hallmark of the Gutzwiller projected Fermi sea, which is
what the Haldane--Shastry ground state is. The extra $\prod_i z_{n_i}$ is a
pure momentum factor: the lattice translation $T: z_n\mapsto z_{n+1}$
multiplies each coordinate by $e^{2\pi i/L}$, and it fixes the total momentum
to the value the ground state actually carries. Dropping it leaves a state that
is \emph{still an exact eigenvector} (residual $\sim10^{-15}$) but in a
different momentum sector: its overlap with the true ground state is zero to
machine precision, at $L=6,8,10$. That is the kind of detail which is invisible
in the algebra and immediate in the numerics, and the reason both variants were
implemented.

The deeper reason such a simple closed form exists is that the $1/r^2$ chain is
not merely integrable but carries a Yangian symmetry, so its spectrum is that
of free "spinons" with no scattering phase --- which is also why (C.3) is a
finite polynomial in $L$ rather than the solution of Bethe equations.

\paragraph{What this implies for a variational circuit.}
Model 2 has a zero-parameter exact ansatz, so running a VQE on it would be an
approximation to something already known. The two honest uses of (C.4) are as a
\emph{warm start} for a perturbed, non-integrable version of the model (say
$1/r^{2+\epsilon}$), where the circuit only has to represent corrections; and
as a benchmark for what happens when an ansatz built on nearest-neighbour bonds
is applied to an all-to-all Hamiltonian. Note also the change in measurement
cost: $H_{\rm HS}$ has $3\binom{L}{2}$ Pauli terms against $3L$ for the XXX
chain ($84$ against $24$ at $L=8$).

\subsection{Model 3: Jordan--Wigner and the SYK model}

\paragraph{\textsf{D.1}--\textsf{D.2} The mapping.}
With $n=N/2$ qubits and $k=0,\dots,n-1$,
\begin{equation}
  \chi_{2k} = \Big(\prod_{j<k}Z_j\Big)X_k,
  \qquad
  \chi_{2k+1} = \Big(\prod_{j<k}Z_j\Big)Y_k .
  \tag{D.0}
\end{equation}
Three cases establish $\{\chi_a,\chi_b\}=2\delta_{ab}$. \emph{(i)} $a=b$: the
$Z$-string squares to the identity and $X_k^2=Y_k^2=\mathbb{1}$, so
$\chi_a^2=\mathbb{1}$. \emph{(ii)} Same qubit, different operator:
$\{X_k,Y_k\}=0$ and the strings cancel. \emph{(iii)} Different qubits $k<l$:
the string of the operator with the larger index contains $Z_k$, and $Z_k$
anticommutes with both $X_k$ and $Y_k$, while every other factor commutes ---
so the pair anticommutes. Verified by explicit matrices for $N=4,\dots,12$.

\paragraph{\textsf{D.3}--\textsf{D.4} Each quartic term is one Hermitian Pauli
string.}
A product of Pauli strings is a Pauli string times a phase in
$\{\pm1,\pm i\}$, so $\chi_i\chi_j\chi_k\chi_l$ is a single string. Reversing a
product of $q$ distinct mutually anticommuting factors costs
$(-1)^{\binom{q}{2}}$; for $q=4$ this is $(-1)^6=+1$, so
\begin{equation}
  \big(\chi_i\chi_j\chi_k\chi_l\big)^\dagger
  = \chi_l\chi_k\chi_j\chi_i
  = \chi_i\chi_j\chi_k\chi_l ,
  \tag{D.4}
\end{equation}
and $H$ is Hermitian for real couplings, with no factor of $i$ needed. (For
$q=2$ the same count gives $(-1)^1=-1$, which is why the free-fermion model of
Section~\ref{sec:warm} carries an explicit $i$.)

\paragraph{\textsf{D.5}--\textsf{D.6} Counting.}
With $J_{ijkl}$ totally antisymmetric, the $4!$ orderings of each index set
contribute equally and cancel the $1/4!$ of the problem statement, so
\begin{equation}
  H_{\rm SYK} = \sum_{i<j<k<l} J_{ijkl}\,\chi_i\chi_j\chi_k\chi_l ,
  \qquad
  \overline{J_{ijkl}}=0,\quad \overline{J^2_{ijkl}}=\frac{6J^2}{N^3},
  \tag{D.5}
\end{equation}
which is exactly $\binom{N}{4}$ Pauli strings on $N/2$ qubits: $N^4/24$ terms
against a Hilbert space of $2^{N/2}$. The Jordan--Wigner $Z$-strings make the
typical string act on $O(n)$ qubits --- measured mean weight $5.1$ out of $8$
at $N=16$ --- so there is no locality to exploit.

\paragraph{\textsf{D.7} Commutation criterion.}
Two Pauli strings either commute or anticommute. On each site they
anticommute iff both factors are non-identity and different; the total sign is
the product over sites, so
\begin{equation}
  [A,B]=0 \iff
  \#\{\,k:\ A_k\neq I,\ B_k\neq I,\ A_k\neq B_k\,\}\ \text{is even}.
  \tag{D.7}
\end{equation}
This is the predicate coloured in the measurement-grouping analysis; at $N=16$
it packs $1820$ strings into $167$ groups, against $549$ under the stricter
qubit-wise criterion.

\paragraph{\textsf{D.8}--\textsf{D.9} The ADAPT gate and its selection
gradient.}
Each pool element is a Pauli string with $P^2=\mathbb{1}$, so
\begin{equation}
  e^{-i\theta P} = \cos\theta\,\mathbb{1} - i\sin\theta\,P ,
  \tag{D.8}
\end{equation}
two sparse matrix--vector products. Appending it at the end of the circuit,
$E(\theta)=\bra{\psi}e^{i\theta P}He^{-i\theta P}\ket{\psi}$ and
\begin{equation}
  \left.\frac{\partial E}{\partial\theta}\right|_{\theta=0}
  = \bra{\psi}i[P,H]\ket{\psi}
  = i\big(\ip{P\psi}{H\psi}-\overline{\ip{P\psi}{H\psi}}\big)
  = -2\,\mathrm{Im}\,\ip{P\psi}{H\psi},
  \tag{D.9}
\end{equation}
using $P^\dagger=P$ and $H^\dagger=H$. This is what ranks the pool; because
each new operator enters at $\theta=0$, the construction is a warm start by
design.

\paragraph{\textsf{D.12} Shot cost.}
For $H=\sum_s c_s P_s$ measured in groups $g$ of mutually commuting strings,
each single shot of group $g$ returns an unbiased estimate of every member
simultaneously, and
\begin{equation}
  \operatorname{Var}\big[\hat E\big]
  = \sum_g \frac{1}{M_g}
    \operatorname{Var}_g\Big[\sum_{s\in g}c_s P_s\Big]
  \;\le\; \frac{\big(\sum_s|c_s|\big)^2}{M},
  \tag{D.12}
\end{equation}
so the worst-case budget for precision $\epsilon$ is
$(\sum_s |c_s|)^2/\epsilon^2$. The implementation samples bitstrings from
$|\psi|^2$ after rotating each qubit into the required basis, so the
correlations within a group are the true ones rather than an assumption; the
estimator is verified unbiased against the exact expectation value.

\paragraph{The reality obstruction, stated precisely.}
A circuit built from $R_Y$ and $CZ$ acting on $\ket{0}^{\otimes n}$ has real
amplitudes in the computational basis, so it can only represent states with
$\mathrm{Im}\,e^{-i\varphi}\psi = 0$ for some global phase $\varphi$. Defining
\begin{equation}
  \kappa(\psi) \;=\; \min_{\varphi}\;
  \frac{\big\|\mathrm{Im}\,e^{-i\varphi}\psi\big\|}{\|\psi\|}
  \;\in\;[0,1],
\end{equation}
a real-amplitude ansatz is fundamentally limited whenever $\kappa>0$. For the
SYK ground state we measure $\kappa = 0.49$ at $N=8$ and $0.66$ at $N=10$,
which is why the $R_Y+CZ$ brickwork returns the \emph{same} energy at every
depth, and why one $R_Z$ per qubit per layer makes it exact at those sizes.

\paragraph{The entanglement obstruction, stated precisely.}
A circuit of depth $d$ built from two-qubit gates can only act across a given
cut $d$ times, so the entanglement entropy it can generate across that cut is
$O(d)$. The SYK ground state has half-cut entropy $76$--$93\%$ of the maximum
$\tfrac n2\ln 2$, growing linearly in $n$; the XXX ground state grows
logarithmically. Hence a constant-depth ansatz cannot scale here --- a
statement about expressivity, independent of the optimiser. At the sizes where
we can diagonalise ($n\le 7$) this is not yet the binding constraint, and
saying otherwise would be overclaiming: it becomes visible as the complex
brickwork degrades from exact at $n=5$ to $8.5\%$ at $n=6$ and $19.4\%$ at
$n=7$.

\subsection{The verification table}

\begin{center}\scriptsize
\input{tables/derivations}
\captionof{table}{Every identity of this appendix, evaluated numerically by
\texttt{verify\_derivations.py}. $47/47$ pass. Entry \textsf{B.6a} is inverted
by construction: it records that the two-term shift rule \emph{fails} for a
four-bond sub-layer, with an observed error of $2.2$.}
\label{tab:derivs}
\end{center}
